\section{Separate record bytes from state publication}
A successful push copies the payload into free ring slots. Once it returns, the
queue owns its own copy; the caller can reuse the source buffer. A successful pop
copies the oldest payload into the destination. The destination is independent of
future ring reuse. This differs from the previous parser's returned slice, which
continued to alias its input.

The demonstration writes the new record before publishing its new tail and used
count. It leaves head unchanged. Pop checks destination capacity before copying,
then publishes its advanced head and reduced used count. It leaves tail unchanged.
Tentative work remains in caller-clobbered registers until the checks pass.
No memory edits are needed for a failed operation.

The descriptor's three state words may be rewritten on success, including a state
word whose value remains unchanged. Their final values must satisfy the operation's
contract. Capacity and the stored buffer pointer must never be written by either
function. A different copy direction or successful state-store order is accepted.
Writes to unrelated slots, even if later restored, are not accepted.

These routines are leaf functions and use temporary registers without altering
saved registers or the stack pointer. A correct stack-using implementation is also
valid if it restores the full calling state. The return boundary supplies the
coherent state observation used by the exercise. Multiple stores do not become a
hardware-atomic transaction, and the example supplies no locks or memory-ordering
protocol for concurrent producers and consumers.
